\newpage
%附录
\appendix
%附录需重新起页，论文附录至少应包括参赛论文的所有源程序代码，如实际使用的软件名称、命令和编写的全部可运行的源程序（含EXCEL、SPSS等软件的交互命令）；通常还应包括自主查阅使用的数据等资料。赛题中提供的数据不要放在附录。如果缺少必要的源程序或程序不能运行（或者运行结果与正文不符），可能会被取消评奖资格。如果确实没有源程序，也应在论文附录中明确说明“本论文没有源程序”。 	

\section{详细图表} % 索引方式	 \ref{xxtb1}  \ref{xxtb2}  \ref{xxtb3}
				  % 详细数据图请参见附录  \ref{xxtb1}。
\subsection*{******详细数据表}%采用不编号格式
\label{xxtb1}

\subsection*{******详细数据表}
\label{xxtb2}

\subsection*{******详细数据表}
\label{xxtb3}


\section{代码程序}
% 插入Python代码
\lstinputlisting[style=Python, caption={Python Example}]{code/code_example.py}

% 插入MATLAB代码
\lstinputlisting[style=MATLAB, caption={MATLAB Example}]{code/code_example.m}



\section{支撑材料}

\begin{enumerate}
	\item 	支撑材料1
	
	\item 	支撑材料2
	
	\item 	支撑材料3
\end{enumerate}

